% PHY 3902 Experiment 3 Lab Report
% Author: Michelle Johnstone, Patrick Geraghty
% Affiliation: University of Ottawa
% Date: October 15, 2025

\documentclass[12pt,a4paper]{article}
\usepackage[utf8]{inputenc}
\usepackage{geometry}
\geometry{margin=1in}
\usepackage{graphicx}
\usepackage{amsmath}
\usepackage{hyperref}
\usepackage{listings}
\usepackage{xcolor}
\usepackage{float}
\usepackage[font=small,labelfont=bf,it]{caption}

% Python and Arduino code formatting
\lstset{
  basicstyle=\ttfamily\footnotesize,
  breaklines=true,
  frame=single,
  backgroundcolor=\color{gray!10},
  keywordstyle=\color{blue},
  commentstyle=\color{green!50!black},
  stringstyle=\color{red},
  numbers=left,
  numberstyle=\tiny\color{gray},
  captionpos=b,
  inputencoding=utf8
}

\begin{document}

% Title
\begin{center}
    {\LARGE\bfseries PHY 3902 Experiment 3 Lab Report}\\[1.5em]
    {\large Author(s): Michelle Johnstone, Patrick Geraghty}\\
    {\large Affiliation: University of Ottawa}\\
    {\large Date: October 15, 2025}
\end{center}

% Abstract
\section*{Abstract}
This experiment analyzed a 650nm laser diode with a built-in series resistor and lens with 5 volts of power applied by shining it on a light-dependent resistor (LDR) used to analyze the intensity of the laser's light. The setup was used for three similar experiments to determine the polarization of the beam as well as observe the effects of a linearly polarized filter in the beam path at various angles. The next step was verifying Brewster's angle by rotating a glass slide in the beam path and analyzing the relationship between the angle of incidence into the slide and the light intensity. The final step was achieving a Gaussian bell curve for the shape and intensity distribution of a beam using a blade to slowly uncover the beam and graphing the gradient of the intensity readings as the blade moved. All the results were as expected; it was determined that the laser's beam was partially linearly polarized and the graph of the light intensity read by the LDR as a function of the angle was sinusoidal with minima when the filter was orthogonal to the direction of the partial polarization of the beam. Brewster's angle was calculated to be 56°, which is approximately the theoretical value of 56.3°, and the intensity distribution bell curve was fairly close to the expected shape with a slight shift to one side.

% Introduction
\section{Introduction}
A linear polarizer limits an electric field incident on it to just one plane. When a beam that is already linearly polarized is incident on a second linearly polarizing filter that is in the same orientation as the first, the transmitted beam will have the same power after both filters. However, if the second filter is placed at an angle $\theta$ relative to the first, the transmitted beam will have intensity $I$ calculated with the equation:
\begin{equation}
I = I_0 \cos^2(\theta)
\end{equation}
This relationship is evident when plotting the intensity as a function of the angle, as the maxima and minima range from 0 to $I$ sinusoidally.

Brewster's angle is the angle of incidence at which, when unpolarized light hits a barrier between two media, the reflected ray is perfectly polarized and the transmitted ray is partially polarized. When p-polarized light is incident on a barrier, the light is entirely transmitted; none of it is reflected. This is due to a phenomenon where the only light reflected from the medium is s-polarized because the reflected and refracted rays are 90° apart. At Brewster's angle, the electric field of the p-polarized light, which oscillates perpendicular to the wave's direction of propagation, would need to oscillate in the direction of the transmitted ray, whereas the electric field of the s-polarized light oscillates perpendicular to the angle of incidence and is reflected. Thus, all of the light is focused into the medium when the incident wave is p-polarized. The determination of Brewster's angle $\theta_B$ is dependent entirely on the indices of refraction of the two media on either side of the barrier. The theoretical value for the expected angle of incidence at which there is no reflected light can be easily calculated using:
\begin{equation}
\theta_{B} = \arctan\left(\frac{n_{2}}{n_{1}}\right)
\end{equation}
where $n_1$ is the index of refraction of the first medium and $n_2$ the index of refraction of the medium the laser beam is propagating into.

For the final portion wherein the Gaussian properties of the laser beam were investigated, a sharp edge is placed on a translational stage in front of a beam and is moved perpendicular to the beam's axis to slowly block more of the beam. As more of the beam is covered, the plot of the intensity distribution is an error function:
\begin{equation}
erf(x) = \frac{2}{\sqrt{\pi}}\int_0^x e^{-t^2} dt
\end{equation}
for which the derivative is the bell curve defined by:
\begin{equation}
f(x) = e^{-x^2}
\end{equation}
This function dictates that the center of the beam is the brightest and the intensity radiates from there and diminishes proportionally to distance from the center.

% Methods
\section{Methods}
For this experiment, the required equipment was few and inexpensive with a fairly simple set up. The laser diode was affixed to a 3D printed stand, and the LDR was set up at equivalent heights with enough separation to allow for the rotating slide to fit in between. The setup was stuck to the workbench using tack and tape in an attempt to stop the components from shifting. A linearly polarizing filter was then attached to a stepper motor on a stand, and the maximum intensity was recorded as the filter was rotated in increments of roughly 0.87° to get a good approximation of the direction of the laser beam's linear polarization to align the laser for the following steps. The first experiment was to analyze polarization by using a stepper motor to slowly turn a linearly polarizing filter in the beam path while measuring the resistance value of the LDR, which is inversely proportional to the intensity. The value of resistance was then graphed as a function of the angle of the filter relative to the starting point. Due to the relationship of the alignment being sinusoidal, using the change in angle of the filter's orientation is an acceptable way of graphing this relation.

The goal for the second section was to investigate Brewster's angle using a 3D printed fitted adaptor to attach a glass microscope slide to a stepper motor so the slide could be moved at increments of 0.18°. The motor was placed so the axis of rotation of the slide was centered on the beam path, so that the incident angle was most accurate and could be rotated 360° to attain 2 readings for Brewster's and achieve the highest accuracy.

The final portion of the experiment involved swapping the glass slide mounted on a motor for a knife-edge style 3D printout on a servo motor. This setup allowed for the beam to be covered and slowly be exposed to the LDR, then graph as a function of the degrees by which the servo was rotated, which is proportional to the amount of the beam uncovered. Using the LDR, it is possible to profile the laser beam, then the gradient of this distribution depicts the Gaussian properties of a beam of light.

% Results and Discussion
\section{Results and Discussion}
For the first portion of the lab that involved checking the polarization of the beam, the results showed that the laser's beam was partially linearly polarized. The graph of the light intensity read by the LDR as a function of the angle was sinusoidal, with minima when the filter was orthogonal to the direction of the partial polarization of the beam. This is illustrated in Figure~\ref{fig:malus}, which shows the measured intensity versus polarizer angle, demonstrating Malus's Law.

\begin{figure}[H]
    \centering
    \includegraphics[width=0.8\textwidth]{../out/malus_plot_double_polarizer.png}
    \caption{Measured intensity versus polarizer angle, demonstrating Malus's Law. Error bars represent measurement uncertainty. The solid line shows the best-fit cosine-squared curve.}
    \label{fig:malus}
\end{figure}

Brewster's angle was calculated to be 56°, which is approximately the theoretical value of 56.3°. Figure~\ref{fig:brewster} shows the measured intensity versus angle of incidence for p-polarized light, with the minimum indicating the Brewster angle, where reflected light is fully polarized.

\begin{figure}[H]
    \centering
    \includegraphics[width=0.8\textwidth]{../out/brewster_plot.png}
    \caption{Measured intensity versus angle of incidence for p-polarized light. The minimum indicates the Brewster angle, where reflected light is fully polarized.}
    \label{fig:brewster}
\end{figure}

The intensity distribution bell curve was fairly close to the expected shape with a slight shift to one side. Figure~\ref{fig:profile} shows the measured intensity profile of the laser beam across its cross-section. The distribution is approximately Gaussian, indicating good beam quality.

\begin{figure}[H]
    \centering
    \includegraphics[width=0.8\textwidth]{../out/laser_profile_plot.png}
    \caption{Measured intensity profile of the laser beam across its cross-section. The distribution is approximately Gaussian, indicating good beam quality.}
    \label{fig:profile}
\end{figure}

Uncertainties were estimated based on the resolution of the stepper motor and the sensitivity of the LDR, with error sources including alignment, ambient light, and detector nonlinearity.

% Conclusion
\section{Conclusion}
The experiment was able to adequately demonstrate the expected graph of the sinusoidal curve created by rotating a linearly polarized filter in front of a linearly polarized beam as well as Brewster's angle of incidence at which a beam polarized parallel to the plane of incidence is fully transmitted with no reflection. Brewster's angle for glass was expected to be 56.3° but was found to be roughly 56° with an error of 0.5%. Finally, the Gaussian bell curve for the incident distribution had the least accuracy as the graph of the gradient of the intensity significantly favoured one side due to a number of potential reasons including oversaturation of the diode and a misalignment of the laser and the LDR.

% References
\section{References}
\begin{itemize}
    \item Hecht, E. (2017). \textit{Optics} (5th ed.). Pearson.
    \item Pedrotti, F. L., Pedrotti, L. S., and Pedrotti, L. M. (2017). \textit{Introduction to Optics} (3rd ed.). Cambridge University Press.
    \item Experimental procedure and code: PHY 3902 course materials and scripts in experiment\_3/src.
    \item Michel Godin; https://uottawa.brightspace.com/content/enforced/523665-2259FS0268421B00/PHY\%203902\_LED\_E\_2023.pdf (Accessed on September 30, 2025)
    \item Dr.~R\"udiger Paschotta (RP); https://www.rp-photonics.com/brewster\_s\_angle.html (Accessed on October 12, 2025)
    \item Vinod Chugani; https://www.datacamp.com/tutorial/gaussian-distribution (Accessed on October 12, 2025)
\end{itemize}

% Supplementary Code
\section{Supplementary Code}
\subsection*{Python Scripts}

    \lstinputlisting[language=Python,caption={malus\_plot.py},label={lst:malus_plot}]{../src/malus_plot.py}
    \lstinputlisting[language=Python,caption={brewster\_plot.py},label={lst:brewster_plot}]{../src/brewster_plot.py}
    \lstinputlisting[language=Python,caption={laser\_profile_plot.py},label={lst:laser_profile_plot}]{../src/laser_profile_plot.py}
    \lstinputlisting[language=Python,caption={error\_propagation.py},label={lst:error_propagation}]{../src/error_propagation.py}

\subsection*{Arduino Sketches}

    \lstinputlisting[language=C,caption={laser\_setup.ino},label={lst:laser_setup}]{../src/laser_setup/laser_setup.ino}
    \lstinputlisting[language=C,caption={laser\_properties\_stepper.ino},label={lst:laser_properties_stepper}]{../src/laser_properties_stepper/laser_properties_stepper.ino}
    \lstinputlisting[language=C,caption={CSVPrinter.h},label={lst:CSVPrinter}]{../src/laser_profile_calibration/CSVPrinter.h}
    \lstinputlisting[language=C,caption={laser\_profile\_calibration.ino},label={lst:laser_profile_calibration}]{../src/laser_profile_calibration/laser_profile_calibration.ino}

% End of document
\end{document}
